% Generated draft: requires mathematical and source review.
\documentclass[11pt]{article}
\usepackage[margin=25mm]{geometry}
\usepackage{fontspec}
\setmainfont{Latin Modern Roman}
\usepackage{amsmath,amssymb,mathtools,unicode-math}
\setmathfont{Latin Modern Math}
\usepackage{enumitem,needspace,xurl,hyperref}
\setlength{\parindent}{0pt}
\setlength{\parskip}{5pt}
\setlength{\emergencystretch}{3em}
\newcommand{\N}{\mathbb{N}}
\newcommand{\Z}{\mathbb{Z}}
\newcommand{\Q}{\mathbb{Q}}
\newcommand{\R}{\mathbb{R}}
\newcommand{\C}{\mathbb{C}}
\newcommand{\F}{\mathbb{F}}
\newcommand{\A}{\mathbb{A}}
\newcommand{\PP}{\mathbb{P}}
\newcommand{\E}{\mathbb{E}}
\newcommand{\eps}{\varepsilon}
\newcommand{\dd}{\,\mathrm{d}}
\begin{document}
\section*{UnsolvedMath ID: 30005488}
\textbf{Source problem code:} catalogue entry
\subsection{Definitions and mathematical statement}
An abelian differential on a smooth compact connected complex curve \(X\) is a nonzero holomorphic one-form \(\omega\) on \(X\). In the usual unprojectivized convention, a stratum of abelian differentials parametrizes isomorphism classes of pairs \((X,\omega)\) with prescribed multiplicities of the zeros of \(\omega\). The supplied record denotes its intended stratum by \(S_\alpha\), but does not specify the genus, the admissible zero data \(\alpha\), whether zeros are labelled, whether a connected component is intended, or whether differentials are considered up to scaling. This description therefore does not fix a convention for the source's \(S_\alpha\).

The source poses two conjectural targets:
\begin{enumerate}
\item Establish the Ax--Schanuel conjecture for \(S_\alpha\). The supplied record does not state that conjecture's maps, ambient spaces, exceptional loci, hypotheses, or dimension inequality, so a precise Ax--Schanuel assertion cannot be recovered from this record alone.
\item Establish the stated Zilber--Pink finiteness assertion. Write
\[
\mathcal A_{\geq 3}(S_\alpha)
:=\{s\in S_\alpha : s\text{ is an arithmetic point of rank at least }3\text{ in the source's sense}\}.
\]
The requested assertion is
\[
\#\mathcal A_{\geq 3}(S_\alpha)<\infty.
\]
Here ``arithmetic point'' and its ``rank'' are unresolved source terms: the supplied record gives neither definition. This displayed formulation records the target but is not a complete intrinsic mathematical statement until those definitions and the convention for \(S_\alpha\) are supplied. The record also leaves the range of strata to which the assertion is intended to apply implicit.
\end{enumerate}

For context, the original excerpt attributes to the works labelled [Ler21] and [KL22] the first steps of a bi-algebraic heuristic for \(S_\alpha\), motivated by its usefulness for diophantine questions on Shimura varieties. It reports a complete geometric characterization of arithmetic points; density of rank-one arithmetic points in the analytification \(S_\alpha^{\mathrm{an}}\); and finiteness of arithmetic points on each algebraically primitive Teichm\"uller curve when the genus satisfies \(g\geq 3\). It further reports linearity of bi-algebraic subvarieties in many cases, without specifying those cases. Finally, contrary to a hope attributed to [KL22], it attributes to [DM] examples of bi-algebraic subvarieties of \(S_\alpha\) that are not linear. These are supplied contextual claims, not independently verified results or additional hypotheses that resolve the missing definitions.
\subsection{Short English statement}
The source asks for an Ax--Schanuel theorem for a stratum of abelian differentials and conjectures that the stratum has only finitely many arithmetic points of rank at least \(3\). The relevant Ax--Schanuel formulation, the meanings of arithmetic point and rank, and the precise stratum convention are missing. Both targets therefore require clarification before a complete mathematical statement can be reviewed.
\subsection{Sources}
\begin{itemize}
\item Local registry: \texttt{data/unsolvedmath/problems.json}, supplied record with permanent UnsolvedMath ID \(30005488\), exact source problem code \texttt{catalogue entry}, title ``Ax--Schanuel and Zilber--Pink for Abelian-Differential Strata.'' The record belongs to the Oberwolfach Reports Open Problems collection and gives proposed year 2023; it identifies no proposer.
\item Supplied report citation: \emph{Komplexe Analysis -- Differential and Algebraic methods in K\"ahler spaces} (2023), Oberwolfach Reports, pp.~965--1030, DOI 10.4171/owr/2023/18. Supplied link: \url{https://doi.org/10.4171/owr/2023/18}. This is the authoritative report citation supplied by the record; its contents have not been independently consulted here.
\item Supplied literature source, labelled ``OWR 2023 report'': \url{https://ems.press/content/serial-article-files/47013}.
\item Supplied literature source, labelled ``Deroin and Matheus'': \url{https://arxiv.org/abs/2303.00642}. The original excerpt instead spells the attribution ``Deroin-Mattheus'' and labels the work [DM]; that correspondence and spelling require source review.
\item The excerpt also cites [Ler21] and [KL22], without full bibliographic information or links. Their authors, titles, and precise statements cannot be recovered from the supplied record.
\item Literature assessment dated exactly \texttt{2026-08-22}: the supplied curation labels the problem open and assesses the Ax--Schanuel and Zilber--Pink directions as conjectural, with existing results covering important special bi-algebraic settings rather than the stated general problem. This is supplied dated curation, not an independently verified assessment; the record itself requests independent verification of time-sensitive status and cited claims.
\item The registry labels the statement \texttt{corrected\_verified} and says that it was checked against the official report and preserves its hypotheses, notation, and target. That verification claim is supplied metadata, not verification performed for this draft.
\end{itemize}
\end{document}
